\documentclass{article}
\usepackage[UTF8]{ctex}
\usepackage{amsmath,amssymb,amsfonts}
\usepackage{enumitem}
\usepackage{geometry}
\geometry{a4paper, margin=1in}
\usepackage{esint}
% === 错题本排版优化（不参与题目标识） ===
\setlength{\parindent}{0pt}
\setlength{\parskip}{0.35em plus 0.10em minus 0.08em}
\setlength{\emergencystretch}{3em}
\setlist{itemsep=0.25em, topsep=0.40em, parsep=0pt}
\allowdisplaybreaks[2]
% === 错题本排版优化结束 ===
\begin{document}
\section*{北京交通大学2026年数学分析真题}
\section*{7．北京交通大学2026年数学分析试卷（回忆版）}
\begin{itemize}

\subsection{题目1}
  \item[1．] 已知 \(a_{1}, b_{1}\) 为正实数 \(\left(a_{1}>b_{1}\right)\) ，且满足
\[
a_{n+1}=\frac{a_{n}+b_{n}}{2}, b_{n+1}=\sqrt{b_{n} a_{n}}(n=1,2, \cdots) .
\]
证明： \(\lim _{n \rightarrow \infty} a_{n}\) 与 \(\lim _{n \rightarrow \infty} b_{n}\) 存在且相等．

\subsection{题目2}
  \item[2．] 假设 \(f(x)\) 在区间 \(I\) 上二阶可导，其中 \(x_{0} \in I\) 为 \(f(x)\) 的极小值点，\(f(x)\) 为 \(I\) 上的严格凸函数，那么 \(f(x)\) 在 \(I\) 上仅有一个唯一的极小值点．
\end{itemize}

\subsection{题目3}
3．证明： \(\lim _{x \rightarrow+\infty} \frac{1}{x} \int_{0}^{x} f(t) \mathrm{d} t=\frac{1}{p} \int_{0}^{p} f(t) \mathrm{d} t\) ，其中 \(f(x)\) 是 \((-\infty,+\infty)\) 上的连续周期函数，其中周期\(p>0\) ．

\subsection{题目4}
4．求不定积分 \(\int \sin ^{4} \mathrm{~d} x\) 与 \(\int \frac{x-5}{x^{2}-3 x+4} \mathrm{~d} x\) ．

\subsection{题目5}
5．设 \(f(x, y)=|x-y| \varphi(x, y)\) ，其中 \(\varphi(x, y)\) 在 \((0,0)\) 的某邻域内连续．

\begin{itemize}
  \item[（1）] \(\varphi(x, y)\) 满足什么条件时，\(f_{x}(0,0), f_{y}(0,0)\) 存在？
  \item[（2）] \(\varphi(x, y)\) 满足什么条件时，\(f(x, y)\) 在 \((0,0)\) 处可微？
\end{itemize}

\subsection{题目6}
6．求过点 \((2,1,3)\) 的平面，使得其与 \(x\) 轴，\(y\) 轴，\(z\) 轴围成的区域的体积为最小值．

\subsection{题目7}
7．求曲面 \(z=\sqrt{3\left(x^{2}+y^{2}\right)}\) 与平面 \(x+y+z=2\) 所围立体的表面积．

\subsection{题目8}
8．\(S\) 为一个封闭曲面，\(n\) 为 \(S\) 的单位外法向量，\(r=(x, y, z)\) ，求证：

\[
\oiint_{S} \frac{\cos (\widehat{n, r})}{|r|^{2}} \mathrm{~d} S= \begin{cases}0, & S \text { 所围区域不包含原点, } \\ 4 \pi, & S \text { 所围区域包含原点. }\end{cases}
\]

\subsection{题目9}
9．设函数列 \(\left\{f_{n}(x)\right\}\) 在 \([a, b]\) 上有连续的导函数，且 \(\left\{f_{n}(x)\right\}\) 在 \([a, b]\) 上收敛于 \(f(x),\left\{f_{n}^{\prime}(x)\right\}\) 在\([a, b]\) 上一致收敛于 \(g(x)\) ，求证：\(f(x)\) 在 \([a, b]\) 上处处可导，且 \(f^{\prime}(x)=g(x)\) ．

\subsection{题目10}
10．求 \(\sum_{n=1}^{\infty} \frac{n^{2}}{(1+a)^{n}}\) ，其中 \(a>0\) ．

\subsection{题目11}
11．\(\left\{f_{n}(x)\right\}\) 一致收敛于 \(f(x)\) ，且每个 \(f_{n}(x)\) 都有界，证明：\(\left\{f_{n}(x)\right\}\) 在定义域中一致有界．

\subsection{题目12}
12．证明：\(\sum_{n=1}^{\infty} n e^{-n x}\) 在 \((0,+\infty)\) 上不一致收敛，但 \(\sum_{n=1}^{\infty} n e^{-n x}\) 在 \((0,+\infty)\) 上连续．
\end{document}
